\documentstyle[%


righttag%


,11pt%


,amssymb%


,russian%               remove in Eng.


,izvru%                 Set izven% in Eng.


]{amsart}





%





\newcommand{\C}{\bold C}


\newcommand{\sgn}{\operatorname{sign}}


\newcommand{\im}{\operatorname{Im}}


\newcommand{\tts}[1]{}





%\hoffset=-15mm%                  For Eng.


\setcounter{page}{37}


\god{2003}


\nomer{6~(493)} %                        Remove in Eng.


%\nomer{Vol.\,47, No.\,2}%               Set in Eng.


%\str{\pageref{begin}--\pageref{end}}%  Set in Eng.


\udk{517.518}





\author{\it Ф.Н.\,Гарифьянов}


% NB:   ^^ remove in Eng.


\title{О проблеме моментов для целых функций экспоненциального типа}


\thanks{Работа выполнена при финансовой поддержке Российского


фонда фундаментальных исследований (код проекта 02-01-00914).}





\date{}


%\pagestyle{myheadings}%         Set in Eng.


\pagestyle{plain} %              Remove in Eng.


\begin{document}


%\thispagestyle{plain}%          Set in Eng.


\maketitle


\label{begin}





{\it Введение.}


В работе исследована проблема моментов


\begin{equation}\label{1}


L (F(t),t^n)\equiv\int_{-\infty}^\infty


\sgn t\exp(-|t|)t^n(F(t)+i^{n+1}F(it))dt=c_n,


\end{equation}


где $c_0, c_1,\dotsc$ --- заданные числа, а $F(z)$ --- целая функция


экспоненциального типа (ц.\,ф.\,э.\,т.) $\sigma$. Ее индикатор должен 


удовлетворять неравенствам $h(\pi k/2)<1$, $k=\overline{0,3}$, т.\,е.


$\sigma\in[0,\sqrt{2})$.


Задача (\ref{1}) рассматривается как обобщение проблем Стильтьеса 


и Гамбургера на случай четырех лучей. Под классом $[\rho, \sigma)$ понимается 


множество целых функций, порядок которых $\lambda \leq \rho$ и тип 


$\mu < \sigma.$ 





В разделе 1 получены три основные теоремы.





\begin{thm}


В классе $[1,\sqrt2/2)$ однородная интерполяционная 


задача $(\ref{1})$ имеет лишь нулевое решение.


\end{thm}





Через $B$ обозначим класс ц.\,ф.\,э.\,т.\ с индикатором


\begin{equation}\label{2}


h(\theta)=(\cos\theta+\sin\theta)/2,\ \ \theta\in [0,\pi/2);


\ \ h(\theta+\pi/2)=h(\theta).


\end{equation}


Пусть их нижние функции имеют на границе сопряженной индикаторной 


диаграммы (единичного квадрата) только логарифмические 


особенности в вершинах $t_1=-t_3=-1/2-i/2$, $t_2=-t_4=1/2-i/2$.


Таким образом, нижние функции голоморфны в плоскости с некоторыми


разрезами, соединяющими точки ветвления $t_j$. Класс соответствующих


нижних функций обозначим через $B_1$.





\begin{thm}


Если $F(z)\in [1;1-\sqrt2/2)$, то $F(z)$ по своим


моментам $\{c_n\}$ восстанавливается рядом


\begin{equation}\label{3}


F(z)=\sum_{k=0}^\infty c_kF_k(z),


\end{equation}


причем 


\begin{equation}\label{4}


F_k(z)\in B,\ \ L(F_k(t),t^n)=\delta_{k,n}.


\end{equation}


\end{thm}





Пусть $D$ --- множество особых точек нижней функции $f(z)$, а 


$\wt{D}$ --- множество особых точек функции 


$(Vf)(z)\equiv f(z-1)+f(z+1)-f(z-i)-f(z+i)$.





\begin{thm}


Ц.\,ф.\,э.\,т.\ $F(z)$, для которой множество 


\begin{equation}\label{6}


A=C\setminus\wt{D}


\end{equation}


связно, единственным образом определяется моментами $(\ref{1})$.





Задача $(\ref{1})$ в классе $B$ $($т.\,е. уже в случае, когда множество $A$


несвязно и $\sigma=\sqrt2/2)$ также всегда имеет единственное решение,


если только


\begin{equation}\label{7}


\underset{n\to\infty}{\varlimsup}{}\sqrt[\uproot{2}n]{|c_n|/n!}<\sqrt2.


\end{equation}


\end{thm}





Если множество $A$ несвязно и $\sigma>\sqrt2/2$, то выделение


классов единственности интерполяционной задачи (\ref{1}) происходит


лишь за счет полного описания множества особых точек $f(z)$. Из


возможных вариантов в статье изучен случай, когда у $f(z)$ имеются


точки ветвления в вершинах $t_j$ и


полюсы в некоторых точках окружности $|z|=\sigma$.





В разделе 2 рассмотрены частные случаи задачи (\ref{1}), сводящиеся к 


классической проблеме моментов Стильтьеса (М.\,С.) --- отыскания 


функции $F(x)$, для которой


\begin{equation}\label{8}


\int_0^\infty F(x)x^n dx=c_n,\ \ n=0,1,\dotsc,


\end{equation}


где $c_0, c_1,\dotsc$ --- заданные числа. Впервые четко


сформулированная в 1894~г. в \cite{Stiltjes}, она 


привлекла


внимание ряда известных математиков (напр., \cite{ShoTam}).


В работах \cite{Hardy1}, \cite{Hardy2} применена теория трансформаций 


Фурье в


комплексной области к исследованию проблемы (\ref{8}). При 


$F(x)e^{\alpha\sqrt x}\in L_1(0,\infty)$, $\alpha>0$,


\tts{[3], [4] are meant} задача


(\ref{8}) сведена к разрешимому в замкнутой форме уравнению


\begin{equation}\label{9}


\int_0^\infty F(x)\cos z\sqrt{x}dx=\sum_{n=0}^\infty


\frac{(-1)^n c_n z^{2n}}{(2n)!},\ \ |\im z|<\alpha.


\end{equation}


Решение единственно с точностью до значений на множествах меры нуль.


Существенно ослабить ограничения на скорость убывания $F(x)$ нельзя,


т.\,к., например, все М.\,С. функции


$$ 


F_\mu(x)=\exp(-


x^\mu\cos\mu\pi)\sin(x^\mu\sin\mu\pi),\ \ 0<\mu<\frac 12,


$$


равны нулю (\cite{Tit48}, гл.\,11, \S\,9). Для функции класса Харди


сумма ряда в уравнении (\ref{9}) должна быть аналитически продолжима


из круга сходимости в полосу $|\im z|<\alpha$ и стремиться к нулю,


когда $z\to\pm\infty$ внутри данной полосы. 





В последнее время вновь возник интерес к проблеме (\ref{8}). В


\cite{Duran} введен класс функций


$$


S_+: F(t)\in C^\infty[0,\infty);\ \ F^{(k)}(0)=0,


\ \ \lim_{t\to\infty}t^mF^{(k)}(t)=0\ \; \forall k,m.


$$


Оказалось, что для любого набора $\{c_n\}\subset\C$ проблема


(\ref{8}) разрешима в классе $S_+$. В \cite{Popov}, в частности,


установлено, что решение $F(x)$ всегда можно выбрать в классе целых


функций фиксированного порядка $\rho>\frac 12$.





Частный случай задачи (\ref{1}) из раздела 1 --- это задача о 


восстановлении целой функции $F(z)\in [1/4;1)$ по М.\,С. функции


$$


%%\begin{equation}\label{10}


\wt{F}(x)=\sqrt{x}F(x)\exp(-\sqrt[4]x).


%%\end{equation}


$$


 


Исследуется и более общая задача: найти целую функцию $F(z)$ с 


помощью условий


\begin{equation}\label{11}


c_n+\lambda\alpha^{5n}\wh{c}_{5n}=b_n,\ \ n=\overline{0,\infty}, 


\ \ |\alpha|<1.


\end{equation}


Здесь $\{\wh{c}_n\}$ --- М.\,С. функции 


$$


%%\begin{equation}\label{12}


\wt{F}(x)=\sqrt{x}F(x)\exp(-\gamma\sqrt[4]x),\ \ \gamma>1,


%%\end{equation}


$$


а для заданной последовательности $\{b_n\}$ выполнено неравенство


(\ref{7}). Подчеркнем принципиальное отличие этой задачи от проблемы


М.С., заключающееся в том, что теперь неизвестен ни один из моментов:


ни моменты $\{с_n\}$ относительно веса


$\exp(- \sqrt[\uproot{2}4]x)$, ни моменты $\{\wh{с}_n\}$ относительно 


веса $\exp(- \gamma\sqrt[\uproot{2}4]x)$. Поэтому здесь


существенно (как и при доказательстве теорем 2 и 3 в случае


несвязного множества $A$) используется разработанная автором теория


особых линейных разностных уравнений с постоянными коэффициентами


\cite{Avtor1}.


\medskip





%% \section{\ }





{\bf 1.} Приступим к исследованию задачи (\ref{1}). Перепишем 


(\ref{1}) как функциональное уравнение


\begin{equation}\label{13}


L(F(t)\exp(-zt),1)=g(z),\ \ z\in D.


\end{equation}


Здесь функция


\begin{equation}\label{14}


g(z)=\sum_{n=0}^\infty \frac{(-1)^n c_nz^n}{n!}


\end{equation}


голоморфна в некоторой окрестности нуля $A_0$. Для нижней функции


$f(z)$ запишем уравнение (\ref{13}) как разностное уравнение


\begin{equation}\label{15}


(Vf)(z)=g(z),\ \ z\in A_0.


\end{equation}


 


Возможны два случая.





а) Множество (\ref{6}) связно (напр., если тип 


$\sigma<\sqrt2/2$).


Тогда $g(z)$ голоморфна на множестве $A$ и $g(\infty)=0$, а


соотношение (\ref{15}) справедливо и при $|z|>1+\sigma$. В таком


случае $g(z)$ ассоциировано по Борелю с ц.\,ф.\,э.\,т. $G(z)$, т.\,е. вместо


(\ref{15}) имеем


\begin{equation}\label{16}


F(z)=G(z)/H(z),\ \ H(z)=\sum_{k=1}^4 (-1)^k\exp(i^kz).


\end{equation}


Ниже будет предложен и другой способ решения частных случаев этой


задачи, не требующий аналитического продолжения суммы ряда (\ref{14})


в окрестность бесконечно удаленной точки.





б) Множество $A$ несвязно. Тогда из (\ref{15}) не следует, что


$(Vf)(z)=g(z)$ при $|z|>1+\sigma$, не говоря уже о том, что $g(z)$


вообще не обязана быть аналитически продолжимой через любую часть


границы $\partial A_0$ или при возможности такого продолжения


удовлетворять условию $g(\infty)=0$. Методы исследования уравнения


(\ref{15}), примененные в п.\,а), уже не работают.





Теперь заметим, что при предположениях теоремы 1 множество (\ref{6})


связно и $G(z)\equiv 0$, т.\,е. с учетом (\ref{16}) имеем $F(z)\equiv


0$, что и завершает ее доказательство.





Покажем, что результаты теоремы 1 неулучшаемы. Пусть $D$ --- единичный


квадрат с вершинами $t_j$, т.\,е. $D\equiv A_0$. Стороны $l_j$


пронумерованы в порядке положительного обхода границы. Пусть


\begin{equation}\label{neprer}


g^+(t)\in C(\overline l_j),\ \ j=\overline{1,4};\ \ g(z)\in A(D).


\end{equation}


В работе \cite{Avtor1} было установлено, что уравнение (\ref{15})


имеет единственное решение, представимое в виде интеграла типа Коши


\begin{equation}\label{coshi}


f(z)=\frac{1}{2\pi i}\int_{\partial D}


\varphi(\tau)(\tau-z)^{-1} d\tau,\ \ z\notin\overline D,


\end{equation}


плотность которого может иметь лишь точки разрыва первого рода в


вершинах. Положим $g(z){\equiv} 1$ и введем функцию


$$


f_0(z) : (Vf_0)(z)=1,\ \ z\in D,\ \Longrightarrow \ (Vf_0^{(k)})(z)=0,


\ \ z\in D,\ \ k\ge1.


$$


Любая система ее производных линейно независима, т.\,е. без фиксации


поведения нижней функции на границе сопряженной индикаторной диаграммы


нельзя гарантировать не только единственности, а даже и конечности


числа решений задачи (\ref{1}). При этом производные


$f_0^{(k)}(z)$ имеют в вершинах $t_j$ неинтегрируемые особенности.


Ц.\,ф.\,э.\,т.


$$


\wt{F}(z)=\int_{|t|=\varepsilon}f'_0(t)\exp(zt)dt,


\ \ \varepsilon >\sqrt2/2,


$$


имеет тип $\sigma=\sqrt2/2$ и нулевые моменты (\ref{1}). Заметим, что


предположение $\sigma<\sqrt2/2$ сразу влечет за собой ситуацию,


описанную в случае~а), откуда $f^\prime_0(z)\equiv 0$, что невозможно.





Для доказательства теоремы 2 введем систему функций 


$$


\{f_m(z) : (Vf_m)(z)=z^m/m!,\ z\in D,\ m=\overline{0,\infty} \}.


$$


Их верхние функции $F_m(z)$ удовлетворяют условию (\ref{4}). Пусть


$D_0$ --- область, ограниченная дугами четырех окружностей


$|z-i^k|=\sqrt2/2$, $k=\overline{1,4}$, и $0\in D_0$. Будем говорить,


что ц.\,ф.\,э.\,т.\ $F(z)$ принадлежит $B(D_0)$, если $f(z)\in A[cD_0]$.


Класс $B(D_0)$ содержит все целые функции класса $[1;\,1- \sqrt2/2)$ и


некоторые из целых функций класса $[1;\sqrt2)$. В работе


(\cite{Avtor2}, \S\,3) была получена формула


$$


f(z)=\sum_{k=0}^\infty c_k f_k(z),\ \ z\in \C\setminus\overline D_0,


$$


причем ряд сходится абсолютно и равномерно на любом компакте


$\overline H\subset\C\setminus\overline D_0$. При этом существенно


использовались равенства


$$


\frac{(-1)^{m+1}}{2\pi i}\int_{\partial D_0}


f_m(z)E^{(k)}(z)dz = \delta_{m,k},


$$


где $E(z)=2z/(z^4-1)$. Подчеркнем, что функции $f_m(z)$ являются


аналитически продолжимыми на множество $D\setminus D_0$ в силу


своего определения. Умножим обе части разложения $f(z)$ в


биортогональный ряд на множитель $\exp(\lambda z)dz$, после чего


проинтегрируем по окружности $L: |z|=\varepsilon$,


$\varepsilon >\sqrt2/2$. В итоге получим искомую формулу (\ref{3}), что


и завершает доказательство. Заметим, что существенно ослабить условия


теоремы~2 нельзя. Достаточно вспомнить построенную выше функцию


$\wt F(z)$, для которой $c_k=0$\; $\forall k$.





Формула (\ref{3}) позволяет обойтись без


аналитического продолжения ряда (\ref{14}) в окрестность бесконечно


удаленной точки.








\begin{pf*}{Доказательство теоремы 3}


Ее первое утверждение --- прямое следствие формулы


(\ref{16}), т.\,к. при $G(z)\equiv 0$ имеем $F(z)\equiv 0$. Второе


утверждение теоремы вытекает из приведенного выше результата 


\cite{Avtor1} о разрешимости разностного уравнения (\ref{15}) в


случае, когда $A_0$ --- квадрат $D$ с вершинами $t_j$. При выполнении


неравенства (\ref{7}) радиус сходимости степенного ряда (\ref{14})


больше, чем $\sqrt2/2$, т.\,е. $g(z)\in A[D]$.


\end{pf*}





\begin{note}


Можно показать, что ц.\,ф.\,э.\,т.\ $F_m(z)$ --- вещественная на $R$,


вполне регулярного роста с индикатором (\ref{2}), причем


$F_m(iz)=-i^mF_m(z)$. Множество ее корней, не лежащих на координатных


осях, может иметь лишь нулевую плотность. Множество ее корней, лежащих


на осях координат ``вдали'' от начала координат, состоит лишь из простых


корней и является $R$-множеством (\cite{Levin}, гл.\,II, \S\,1). Функция


$F_m(z)$ --- каноническое произведение Вейерштрасса нулевого рода и


$\exists\gamma :|F_m(z)|\le\gamma\exp(\sqrt2|z|/2)/m!$\; $\forall m$.


Такие результаты следуют из свойств нижних функций $f_m(z)$,


установленных в (\cite{Avtor1}, \S\,3).


\end{note}





\begin{note}


Кроме ряда (\ref{3}), равенство (\ref{4}) порождает


еще один биортогональный ряд


$P(z)=\sum\limits_{m=0}^\infty a_mz^m$,


т.\,е. обычный ряд Маклорена. Легко установить, что если $P(z)$ --- целая


функция класса $[1;1/2)$, то ее коэффициенты $a_m$ могут быть


вычислены с помощью (\ref{4}). Это позволяет получить


некоторую информацию о свойствах ц.\,ф.\,э.\,т.\ $F_m(z)$ (см. раздел 2).


\end{note}





Выделим некоторые классы ц.\,ф.\,э.\,т., в которых проблема (\ref{1}) 


имеет единственное решение.





1) При $\sigma =\sqrt2/2$ в качестве такого класса выберем $B$. Условия


(\ref{neprer}) заведомо выполнены, если справедливо неравенство


(\ref{7}).





2) При $\sigma\in(\sqrt2/2,\sqrt{2})$ берем класс ц.\,ф.\,э.\,т.\ $B_2$.


Поведение


ассоциированных с ними по Борелю нижних функций $f(z)$ на $\partial D$


идентично поведению функций из класса $B_1$. Но у $f(z)$ есть еще


особенность --- простой полюс в точке $z_0=\sigma\exp(i\varphi)$ с


фиксированным вычетом $\beta$. Аргумент $\varphi$ выбран так, чтобы


$(V\wt{f})(z)\in A[D]$, где $\wt{f}(z)=\beta(z-z_0)^{-1}$.


Другими словами, точки $z_0$ --- это те точки окружности $|z|=\sigma$,


которые лежат вне замыкания множества $\wt{D}$. Уравнение


$$


(Vf)(z)=-(V\wt{f})(z)+g(z),\ \ z\in D,


$$


имеет единственное решение $f(z)\in B_1$. Решением задачи (\ref{1})


будет ц.\,ф.\,э.\,т., ассоциированная по Борелю с нижней функцией


$f(z)+\wt{f}(z)$.





Перечисленные классы единственности не исчерпывают всех возможных


вариантов. Здесь прослеживается связь с теорией краевых задач для


аналитических функций, когда индекс задачи, от которого зависит число решений


или условий разрешимости, четко определяется выбранным классом


граничных значений.





%% \section{\ }





{\bf 2.} Подробнее изучим свойства функций $F_{4m+3}(z)$. Имеем


\begin{equation}\label{18}


F_{4m+3}(z)=z^2P_m(z^4),


\ \ \int_0^\infty t^{k+1/2}P_m(t)\exp(-\sqrt[\uproot{2}4]t)dt=\delta_{m,k},


\end{equation}


где $P_m(t)$ --- вещественная на действительной оси функция с 


тригонометрическим индикатором


$h(\theta)=[\cos(\theta/4)+\sin(\theta/4)]$. Введем функции


\begin{equation}\label{19}


P(z)=\sum_{j=1}^n\wt{P}_j(z)P_{m_j}(z),


\end{equation}


где $\wt{P}_j(z)$ --- вещественные на действительной оси целые функции класса


$[1/4;\varepsilon]$, $\varepsilon<1/2$. Пусть $\{a_k\}_0^\infty$ ---


множество положительных корней нечетной кратности целой функции


(\ref{19}). Образуем каноническое произведение Вейерштрасса


$\Theta(z)=\prod\limits_{k=1}^\infty\big(1-\frac{z}{a_k}\big)$.


Подчеркнем, что все корни $\Theta(z)$ простые. Докажем, что у


канонического произведения $\Theta(z)$ тип $\sigma\ge1/2-\varepsilon$.


Предположим противное. Тогда существует интеграл


$$


J=\int_0^\infty x^{k+1/2}P(x)\Theta(x)\exp(-\sqrt[\uproot{2}4]x)dx


$$


и $J\ne 0$ в силу знакопостоянства произведения $P(x)\Theta(x)$. С другой


стороны, почленное интегрирование ряда Маклорена целой функции


$x^k\wt{P}_j(x)\Theta(x)$ типа $\sigma<1/2$ с учетом (\ref{18})


приводит к тому, что $J=0$ при $k>\max\limits_j m_j$ (см. замечание~2).





Изложим ход решения уравнения


$$


(Vf_1)(z)+\lambda(f_1[\mu(z)+\gamma]+f_1[\mu(z)-\gamma]


-f_1[\mu(z)+\gamma i]-f_1[\mu(z)-\gamma i])=g(z),\ \ z\in D,


$$


где $\mu(z)\in A[D]$, $\mu(iz)=i\mu(z)$, $\mu(\overline D)\subset D$,


$\gamma>1$, аналогичный проведенному в


\cite{Avtor1} для $\lambda=0$. Используем интегральное представление


(\ref{coshi}), неизвестная плотность которого удовлетворяет условию


$$


\varphi(t)+\Theta_t\varphi(\alpha(t))=0.


$$


Здесь $\Theta_t=\{1,\, t\in l_1\cup l_3;\, -1,\, t\in l_2\cup l_4\}$, а


$l_1$ --- сторона квадрата с концами в точках $t_1$ и~$t_2$. Сдвиг


$\alpha(t)$ на стороне $l_k$ определен формулой $\alpha(t)=t+i^k$,


$k=\overline{1,4}$, т.\,е.\ $\alpha(t): \partial D\to \partial D$ с


изменением ориентации. Поэтому $\alpha(\alpha(t))\equiv t$,


$t\notin\{t_k\}$. Теперь разностное уравнение сводится к


функциональному уравнению


\begin{equation}\label{20}


\Phi(z)\equiv\frac{1}{2\pi i}\int_{\partial D}\varphi(\tau)


A(z,\tau)d\tau =g(z),\ \ z\in D,


\end{equation}


с ядром $A(z,\tau)=E(\tau-z)+\gamma^{-1}E((\tau-\mu(z))/\gamma)$.


С помощью аналога формулы Сохоцкого 


$$


\Phi^+(t)=-\frac{\Theta_t}{2}\varphi(\alpha(t))+\Phi(t)


$$


перейдем к интегральному уравнению Фредгольма


\begin{equation}\label{21}


(T\varphi)(t)\equiv\varphi(t){+}\frac{1}{2\pi i}\int_{\partial D}


(A(t,\tau)-\Theta_t\Theta_\tau A(\alpha(t),\alpha(\tau)))\varphi(\tau) 


d\tau = g^+(t)-\Theta_tg^+(\alpha(t)),\ \ t\in\partial D.


\end{equation}


Заметим, что $(T\varphi)(t)\equiv\Phi^+(t)-\Theta_t\Phi^+(\alpha(t))$.


Это позволяет с помощью теории краевой задачи Карлемана установить


эквивалентность функционального уравнения (\ref{20}) и интегрального


уравнения (\ref{21}). Ядро уравнения (\ref{21}) ограничено на


$\partial D$. Оператор $T$ удобно рассматривать на банаховом


пространстве $\wt C(\partial D)$ --- множестве функций


$\varphi(t)\in C(\overline{l}_j)$, $j=\overline{1,4}$, с нормой


$\|\varphi\|=\max|\varphi(t)|$, $t\in\partial D$. Очевидно,


$T:\wt C(\partial D)\to\wt C(\partial D)$. При $\lambda=0$


обратимость оператора $T$ была установлена в \cite{Avtor1} применением


принципа сжимающих отображений. Этот результат остается


справедливым и при достаточно малых $|\lambda|\ne0$. В частности, при


$\mu(z)=\alpha z^5$, $g(iz)=-ig(z)$, $|\alpha|<1$, для функций


$F(x)=F_1(\sqrt[4]x)\sqrt{x}$, где $F_1(x)$ ассоциирована по Борелю


с $f_1(x)$, получим задачу (\ref{11}).





\begin{note}


Предложенные методы позволяют свести проблему моментов


$$


\int_0^\infty F(x)e^{-x}x^k dx


+\lambda\int_0^{-\infty}F(x)e^{-|x|}x^k dx+


\beta\int_0^{i\infty} F(x)e^{-|x|}x^k dx+


\lambda\beta\int_0^{-i\infty} F(x)e^{-|x|}x^k dx=c_k


$$


в классе функций с фиксированным индикатором (\ref{2}) к исследованию


некоторого интегрального уравнения Фредгольма 2-го рода.


\end{note}





Пусть $D$ --- фундаментальный многоугольник конечно-порожденной


разрывной группы дробно-линейных преобразований, ограниченный конечным


числом дуг окружностей, причем $D$ выбран так, что среди его вершин


нет предельных точек группы. Считаем, что любая вершина $D$ является


общей для четного или бесконечного числа конгруэнтных фундаментальных 


многоугольников.


Рассмотрим функциональные уравнения (\ref{20}), где


\begin{equation}\label{22}


A(z,\tau)=\sum^\infty_{k=0}\frac{1}{\tau-\sigma_k(z)}.


\end{equation}


Обозначим через $S_r(z)$ порождающие преобразования группы, 


для которых конгруэнтные многоугольники $S_r(\overline{D})$ имеют с


$\overline{D}$ общую сторону. Пусть $m(j)$ --- наименьшее число


порождающих преобразований $S_r(z)$ и преобразований, обратных к ним,


суперпозицией которых является преобразование $\sigma_j(z)$. В этой


суперпозиции каждое преобразование считается столько раз, какова его


кратность. При этом на разных местах могут быть одни и те же


преобразования $S_r(z)$ или $S_r^{-1}(z)$, поскольку дробно-линейные


функции, вообще говоря, не коммутируют. Введем такую функцию


$\alpha(t)$, что $\alpha(\alpha(t))=t$, а на сторонах $l_k$ она


совпадает с порождающими преобразованиями группы


$\alpha(l_k)=l'_k$. Здесь $l_k$ и $l^\prime_k$ --- конгруэнтные


стороны границы $\partial D$.





Пусть в формуле (\ref{22}) сумма берется по таким преобразованиям


группы, для которых числа $m(k)$ нечетны и, кроме того,


$\sigma_k(\overline D)\cap \overline D \ne\emptyset$ (свойства таких


лакунарных ядер см. в \cite {Avtor3}). Тогда


функциональное уравнение (\ref{20}), с одной стороны, равносильно


проблеме моментов


$$


\sum_{k=0}^\infty\int_{\arg \tau=\varphi_k}F(\tau)\frac{d^j}{dz^j}


\exp(-\sigma_k(z)\tau)|_{z=a}d\tau=g^{(j)}(a),\ \ a\in D,


\ \ j=\overline{0,\infty}.


$$





С другой стороны, оказывается возможным с помощью формул Ю.В.\,Сохоцкого


свести его к исследованию интегрального уравнения Фредгольма второго


рода с ядром


$$


K(t,\tau)=A(t,\tau)-\alpha^\prime(\tau)A(\alpha(t),\alpha(\tau)),


$$


ограниченным на $\partial D$.





\begin{thebibliography}{99}





\bibitem{Stiltjes}


Стильтьес Т.И. {\em Исследование о непрерывных дробях}. --


Харьков--Киев: ГНТИ Украины, 1936. -- 156 c.


\bibitem{ShoTam}


Ахиезер Н.И. {\em Классическая проблема моментов и некоторые вопросы


анализа, связанные с нею}. -- М.: Физматгиз, 1961. -- 301 с.


\bibitem{Hardy1}


Hardy G.H. {\em On Stieltjes ``probleme des


moments''} // Messenger of Math. -- 1916. -- V.\,46. -- P.\,175--182.


\bibitem{Hardy2}


Hardy G.H. {\em On Stieltjes ``probleme des


moments''} // Messenger of Math. -- 1917. -- V.\,47. -- P.\,81--88.


\bibitem{Tit48}


Титчмарш Е. {\em Введение в теорию интегралов Фурье}. -- М.--Л.:


ГИТТЛ, 1948. -- 479 с.


\bibitem{Duran}


Duran A.I. {\em The Stieltjes moments problem for rapidly


decreasing functions} // Proc. Amer. Math. Soc. -- 1989. -- V.\,107.


-- P.\,631--641.


\bibitem{Popov}


Попов А.Ю. {\em О проблеме моментов быстро убывающих


функций} // Матем. заметки. -- 1996. -- Т.\,60. -- \No\,1. -- C.\,66--74.


\bibitem{Avtor1}


Гарифьянов Ф.Н. {\em Проблема обращения особого интеграла


и разностные уравнения для функций, аналитических вне квадрата} // Изв.


вузов. Математика. -- 1993. -- \No\,7. -- C.\,7--16.


\bibitem{Avtor2} 


Гарифьянов Ф.Н. {\em Трансформации биортогональных систем


и некоторые их приложения}. II // Изв. вузов. Математика.


-- 1996. -- \No\,8. -- C.\,13--25.


\bibitem{Levin}


Левин Б.Я. {\em Распределение корней целых функций}. -- М.:


ГИТТЛ, 1956. -- 632 с.


\bibitem{Avtor3} 


Аксентьева Е.П., Гарифьянов Ф.Н. {\em О лакунарных аналогах тэта-ряда


Пуанкаре и их приложении} // Сиб. матем. журн. -- 2002.


-- Т.\,43. -- \No\,5. -- C.\,977--986.





\end{thebibliography}





\vskip 2cm





\post{\it Казанский государственный}{\it Поступила}


\post{\it энергетический университет}{18.05.2001}





\label{end}


\end{document}


